\documentclass[10pt,a4paper]{amsart}
\usepackage[margin=19mm]{geometry}
\usepackage{amsmath,amssymb,mathtools}
\usepackage[scr=rsfs,cal=euler]{mathalfa}
\usepackage{xfrac}
\usepackage[unicode,hidelinks,bookmarks=false]{hyperref}
\newcommand{\ie}{i.e.\ }
\newcommand{\cf}{cf.\ }
\newcommand{\resp}{respectively\ }
\newcommand{\Subsection}{\subsection}
\providecommand{\nobreakdash}{\nobreak\hbox{-}}
\setlength{\emergencystretch}{2em}
\multlinegap0pt
\theoremstyle{plain}
\newtheorem{theorem}{Theorem}
\newtheorem{lemma}{Lemma}
\theoremstyle{definition}
\newtheorem{definition}{Definition}
\DeclareMathOperator{\Sym}{\mathrm{Sym}}
\title[Generating the symmetric group by three prefix reversals]{Generating the symmetric group by three prefix reversals: the case m = n - 2 of the three-prefix-reversal classification (printed as Theorem 4)}
\author{Sa\'ul A. Blanco, Mikhail P. Golubyatnikov, Elena V. Konstantinova, Natalia V. Maslova, Luka A. Nikiforov}
\date{Source: arXiv:2511.16959v3 (CC BY 4.0) (CC BY 4.0)}
\begin{document}
\maketitle

\noindent\emph{Source and licence.} Generating the symmetric group by three prefix reversals. Version arXiv:2511.16959v3 (CC BY 4.0), distributed under the Creative Commons Attribution 4.0 International licence, http://creativecommons.org/licenses/by/4.0/, as recorded in the arXiv licence metadata for this exact version and shown on the saved abstract page https://arxiv.org/abs/2511.16959v3. Full licence notice: licenses/arxiv-2511.16959-CC-BY-4.0.txt.\par\smallskip

\noindent\textit{Editorial scope.} Counted unit: the case m = n - 2 of the classification of three-element sets of prefix reversals that generate the symmetric group, printed in the article as Theorem 4 (source0.tex lines 663--665, one \texttt{theorem} environment with source label \texttt{thm:n-2}), delivered together with the article's complete original proof of it (source0.tex lines 667--837, one \texttt{proof} environment of 171 lines immediately adjacent to the statement, with no gap and no deferral). The proof is the article's longest case analysis: it reduces the cases k in {2,3} to the two earlier theorems, then, for the parity-compatible cases, constructs an (n-2)-cycle explicitly, multiplies it by the third prefix reversal and reads off from the resulting permutation chains a transposition that swaps two adjacent entries of that cycle, after which the article's classical lemma on an n-cycle and a transposition gives generation; for the remaining parity-compatible case it conjugates the transposition (1 2) by powers of the (n-2)-cycle to obtain all transpositions between odd and even entries, that is, the Coxeter generators; and for equal parities it uses the parity obstruction recorded in the introduction. No mathematics is written, repaired or re-derived: statement and proof are the article's own text line for line, in the article's own notation, and no line of either is omitted, substituted or re-flowed.

Required context retained uncounted: the article's definition of the prefix reversals with its labelled display (source0.tex lines 136--143), which fixes the notation $r_i$ and the labelled equation the article refers to; the five necessary conditions on a generating triple (source0.tex lines 150--156), which the proof's final step invokes as Observations 3 and 5; and the article's classical lemma on an n-cycle and a transposition, with its label (source0.tex lines 277--278), which the proof invokes directly.

One declared omission, and nothing else. The article attributes that classical lemma to a published reference by a citation command inside the lemma's header. The article's own bibliography is a BibTeX database that is not part of the delivered source, and this package therefore carries no bibliography block and no bibliography entry; delivering a citation command without its entry would leave an unresolved reference, so the attribution header was omitted from the delivered line under the declared omission receipt carried in delivery-evidence.json. Exactly one string was removed, on exactly one source line, and the receipt names it literally; the mathematical content of the lemma is untouched. No other citation command occurs in any retained line.

Editorial formatting only, in addition: an AMS \texttt{amsart} document that restates the macro definitions the retained lines use (the article's own $\Sym$ operator) and the article's own theorem-environment declarations. The article's own publisher class file \texttt{dmtcs-episciences.cls} and its BibTeX database are neither shipped in this package nor used to build it. Consequently the retained environments print here in the order in which they occur, so the retained lemma prints as Lemma 1 and the counted theorem as Theorem 1, whereas the article prints them as Lemma 2.4 and Theorem 4; the article's per-section numbering is not reproduced, the equations of the retained spans are numbered anew in order of appearance, and every cross-reference inside the retained text resolves to this wrapper's numbering. That is the only change to the numbering. No figure, no image-inclusion line and no drawing code occur in any retained line, so no artwork is delivered or needed, although the article contains figures elsewhere.

Not retained: the article's abstract, its other cases, its conjectures and computational results and its bibliography; none of that material is used as a step of the delivered proof, which is complete as the author wrote it. The package ships the article's concatenated source verbatim as \texttt{sources/189541/source0.tex}, and every delivered excerpt is an exact line range of that file. No mathematical statement, hypothesis, estimate or conclusion has been rewritten, repaired or added, and printed slips, if any, are preserved literally.
The symmetric group $\mathrm{Sym}_n$ is also generated by a set of $n-1$ involutions known as \emph{prefix reversals}, defined by
\begin{equation} \label{e1}
r_i=\begin{pmatrix} 
1 & 2 & \ldots & i-1 & i & i+1 & \ldots & n \\ 
i & i-1 & \ldots & 2 & 1 & i+1 & \ldots & n
\end{pmatrix},
\end{equation}
where $1<i\leqslant n$. The action of $r_i$ is to reverse the order of the first $i$ elements while leaving the remaining $n-i$ elements fixed. Equivalently, $r_i$ is the product of $\lfloor i/2 \rfloor$ disjoint transpositions, $r_i=(1\ i)(2\ (i-1))\cdots(\lfloor i/2\rfloor\ (i+1-\lfloor i/2\rfloor))$.

\begin{enumerate}
    \item[$(1)$] One of the prefix reversals must be $r_n$, since otherwise no element can be moved to or from position $n$. 
    \item[$(2)$] One of the prefix reversals, say $r_i$, must satisfy $\lfloor n/2 \rfloor < i < n$, since otherwise if we take $g\in \langle r_n, r_j, r_i\rangle$ with $i<j \leqslant n/2$, and if $u <v \leqslant j$ then either $g(u) \leqslant n/2$ and $g(v) \leqslant n/2$ or $g(u) > n/2$ and $g(v)> n/2$. Thus, in this case, $2$-transitivity is lost, while $\Sym_n$ is known to be $2$-transitive in its natural action for each $n$. Moreover, if $n$ is odd, then $i>\lceil n/2\rceil$ since otherwise $\lceil n/2\rceil$ will remain fixed. 
    \item[$(3)$] At least one prefix reversal must have an even index, since otherwise no element can be moved from an odd-numbered position to an even-numbered position. 
    \item[$(4)$] At least one prefix reversal must be an odd permutation, since otherwise the generated group consists only of even permutations. 
    \item[$(5)$] The indices must satisfy the condition $\gcd(i,j,k)=1$. Otherwise, if $\gcd(i,j,k)=\ell>1$ then the elements of $\{1,2,\ldots,n\}$ are partitioned into blocks of $\ell$ consecutive positions that can never be separated. In particular, at least one index should be odd.
\end{enumerate}

\begin{lemma}\label{lem:janusz} Let $\sigma$ be an $n$-cycle and $\tau=(a \ b)$ be a transposition. Let $q$ be an integer satisfying $\sigma^q(a)=b$. Then $\langle \sigma,\tau\rangle=\Sym_n$ if and only if $\gcd(q,n)=1$.
\end{lemma}

\begin{theorem}\label{thm:n-2}
Let $H=\langle r_n, r_{n-2}, r_k\rangle$, with $2\leqslant k < n-2$ and $n\geqslant 5$. Then $H=\Sym_n$ if and only if $n$ and $k$ are of different parities. 
\end{theorem}

\begin{proof} The cases $k \in \{2, 3\}$ were already handled in Theorems~1 and~3, respectively. Therefore, one can assume that $k>3$. We first establish the sufficient statement. Suppose $n$ is even and $k$ is odd. Then $H$ contains the following $n$-cycle:
$$\sigma = r_{n-2}r_{n}=(n\ n-2\ n-4\dots 2\ n-1\ n-3 \dots \ 3 \  1).$$ 
To see this, observe that the following holds:
\[
\sigma=r_{n-2}r_n(i)=
\begin{cases}
	n, & \text{if }i=1;\\
	n-1, & \text{if }i=2;\\
	i-2, & \text{if }3\leqslant i\leqslant n.\\
\end{cases}
\]
Starting with \(n\) we obtain the chain:
\[
n\rightarrow n-2\rightarrow n-4 \rightarrow\cdots\rightarrow 2.
\]
Since \(2\rightarrow n-1\) the cycle can be continued as follows:
\[
n-1\rightarrow n-3\rightarrow n-5\rightarrow 3 \rightarrow\cdots \rightarrow 1.
\]
Finally, we have \(1\rightarrow n\) which gives the sought $n$-cycle $\sigma$.

Since $k$ is odd and $k > 3$ then for $k=2m+3$ let us consider the following permutation:
\[
\rho  = (r_{n-2}r_n)^{\frac{k-3}{2}}r_k = \sigma^{m}r_k.
\] 
For $i \geqslant 2m+1=k-2$, applying $\sigma^{m}$ simply subtracts $2$, $m$ times; thus, we have $\rho(1) = \sigma^{m}(2m + 3) = 3$, $\rho(2) = \sigma^{m}(2m + 2) = 2$ and $\rho(3) = \sigma^{m}(2m + 1) = 1$.
Therefore, the restriction of $\rho$ to the set $\{1,2,3\}$ satisfies the following: $$\rho\vert_{\{1,2,3\}} = (1\ 3).$$

Let $4 \leqslant i \leqslant n$. If $i \leqslant k$ then we have: 
\[
\rho(i) = \sigma^{m}r_k(i) = \sigma^{m}(k+1-i) = \sigma^{m}(2m+4-i),
\]
%Note that for $x\leqslant m$ we have: 
%\[
%\sigma^{m}(2x) = n-1-2(m-x)
%\quad\text{ and  }\quad 
%\sigma^{m}(2x+1) = n-2(m-x-1).
%\]
and
\[
\sigma^{m}(2x) = n-1-2(m-x), \; \text{for } 1 \leqslant x \leqslant m;
\]
\[
\sigma^{m}(2x+1) = n-2(m-x-1), \; \text{for } 0 \leqslant x \leqslant m-1.
\]

If $i$ is even then $2m+4-i$ is even and the following holds:
\[
\sigma^{m}(2m+4-i) = n-1 - 2m + 2m+4-i = n + 3 - i.
\]

If $i$ is odd then $2m+4-i$ is odd and the following holds:
\[
\sigma^{m}(2m+4-i) = n + 5 - i.     
\]

If $i > k = 2m+3$ then $\rho(i) = \sigma^{m}(i) = i - 2m$. Thus, we have:
\[
\rho(i)=
\begin{cases}
	3, & \text{if }i=1;\\
	2, & \text{if }i=2;\\
	1, & \text{if }i=3;\\
	n+3 - i, & \text{if }4\leqslant i\leqslant k \mbox{ and } i \text{ even};\\
	n+5 - i, & \text{if }4\leqslant i\leqslant k \mbox{ and } i \text{ odd};\\
	i - 2m, & \text{if }k < i\leqslant n.\\
\end{cases}
\]

Starting with \(4\) we obtain the chain:
\[
4 \rightarrow n-1 \rightarrow n-1-2m \rightarrow n-1-4m \rightarrow\cdots\rightarrow a_1, 
\]
where $a_1 \in \{4, \dots, k\}$ and $a_1 \equiv n-1 \pmod{2m}$. Since $n-1$ is odd then $a_1$ is odd and we have: 
\[
a_1 \rightarrow n + 5-a_1 \rightarrow n + 5-a_1 - 2m \rightarrow\cdots\rightarrow a_2,
\]
where $a_2 \in \{4, \dots, k\}$ and  $a_2 \equiv n + 5-a_1 \equiv  n + 5 -(n-1) \equiv 6 \pmod{2m}$. Thus, we have the following chain:
\[
4 \rightarrow \cdots \rightarrow 6 \rightarrow \cdots \rightarrow 8 \rightarrow \cdots \rightarrow k-1.
\]
Since $k-1 + 2 \equiv 2m+4 \equiv 4 \pmod{2m}$ we continue as follows: 
\[
k-1 \rightarrow \cdots \rightarrow 4.
\]
It remains to observe that each interval $a \rightarrow \cdots \rightarrow a+2$ contains exactly $2m$ terms and all of them belong to the same residue class modulo $2m$, namely, the class containing $n+3-a$. Thus, all odd values $\{5,7,\dots,k\}$ occur in the sequence, and consequently $\rho|_{\{4,5,\dots,n\}}$
is an $(n-3)$-cycle.

Hence, the following two equations hold: 
\begin{align*} (r_{n-2}r_n)^{\frac{k-3}{2}}r_k &= (1\ 3) \underbrace{(\dots \ldots \ldots \ldots \ldots)}_{\text{$(n-3)$-cycle}}, \\
((r_{n-2}r_n)^{\frac{k-3}{2}}r_k)^{n-3}&= (1\ 3). 
\end{align*}
Therefore, having the $n$-cycle $r_{n-2}r_n$ and the transposition $(1\ 3)$ that swaps two adjacent entries in this cycle, by Lemma~\ref{lem:janusz}, we have $H=\Sym_n$.

Now, let us assume that $n$ is odd and $k$ is even. Then we have: 
$$\sigma=r_{n-2} r_n=(n\ n-2\ n-4\dots 1)(n-1\ n-3 \dots 2).$$ 
%%%%%%%%%%%%%% Misha's proof. Feel free to moderate.
Indeed, since $n$ is odd we have the following two chains:
\[
n \rightarrow n-2 \rightarrow n-4 \rightarrow\cdots\rightarrow 1 \rightarrow n
\]
and 
\[
n-1 \rightarrow n-3 \rightarrow n-5 \rightarrow\cdots\rightarrow 2 \rightarrow n-1.
\]

Since $k$ is even and $k > 2$, we can set $k=2m+2$ and consider the following permutation:
\[
\rho  = (r_{n-2}r_n)^{\frac{k-2}{2}}r_k = \sigma^{m}r_k
\]
with $\rho(1) = \sigma^{m}(2m+2) = 2$ and $\rho(2) = \sigma^{m}(2m+1) = 1$ which gives us: $$\rho\vert_{\{1,2\}} = (1\ 2).$$

Let $3 \leqslant i \leqslant n$. If $i \leqslant k$ then the following hold: 
\[
\rho(i) = \sigma^{m}r_k(i) = \sigma^{m}(k+1-i) = \sigma^{m}(2m+3-i) 
\]
and we have: 
\[
\sigma^{m}(2x) = n-1-2(m-x), \; \text{for } 1 \leqslant x \leqslant m
\]
and
\[
\sigma^{m}(2x+1) = n-2(m-x-1), \; \text{for } 0 \leqslant x \leqslant m-1.
\]

If $i$ is odd then $2m+3-i$ is even and we have: 
\[
\sigma^{m}(2m+3-i) = n-1 - 2m + (2m+3-i) = n + 2 - i.
\]
If $i$ is even then $2m+3-i$ is odd and we have: 
\[
\sigma^{m}(2m+4-i) = n + 4 - i.     
\]

If $i > k = 2m+2$ then $\rho(i) = \sigma^{m}(i) = i - 2m$. Hence, we have:
\[
\rho(i)=
\begin{cases}
	2, & \text{if }i=1;\\
	1, & \text{if }i=2;\\
	n+4 - i, & \text{if }3\leqslant i\leqslant k \mbox{ and } i \text{ even};\\
	n+2 - i, & \text{if }3\leqslant i\leqslant k \mbox{ and } i \text{ odd};\\
	i - 2m, & \text{if }k < i\leqslant n.\\
\end{cases}
\]

It remains to show that the restriction of $\rho$ to $\{3,4,\dots,n\}$ is a single $(n-2)$-cycle. Indeed, starting with $3$ we have the chain as follows:
\[
3 \rightarrow n-1 \rightarrow n-1-2m \rightarrow n-1-4m \rightarrow \cdots \rightarrow b_1, 
\]
where $b_1 \in \{3, \dots, k\}$ and $b_1 \equiv n-1 \pmod{2m}$. Since $n-1$ is even, we have $b_1$ is even and  
\[
b_1 \rightarrow n+4-b_1 \rightarrow n+4-b_1-2m \rightarrow n+4-b_1-4m \rightarrow \cdots \rightarrow b_2, 
\] 
where $b_2 \in \{3, \dots, k\}$ and $b_2 \equiv n + 4-b_1 \equiv  n + 4 -(n-1) \equiv 5 \pmod{2m}$. Thus, finally we have the following chain:
\[
3 \rightarrow \cdots \rightarrow 5 \rightarrow \cdots \rightarrow 7 \rightarrow \cdots \rightarrow k-1 \rightarrow 3.
\]

It remains to see that each interval $a \rightarrow \cdots \rightarrow a+2$ contains exactly $2m$ terms and all of them belong to the same residue class modulo $2m$, namely, the class containing $n+2-a$. Thus, all even values $\{4,6,\dots,k\}$ occur in the sequence, and consequently $\rho|_{\{3,4,\dots,n\}}$ is an $(n-2)$-cycle.

Thus, finally the following two equations hold: 
\begin{align*}
(r_{n-2}r_{n})^{\frac{k-2}{2}}r_k &= (1\ 2) \underbrace{(\dots \ldots \ldots \ldots \ldots)}_{(n-2)\text{-cycle}}, \\
((r_{n-2}r_{n})^{\frac{k-2}{2}}r_k)^{n-2} &=(1\ 2).
\end{align*}

Since the lengths of the two cycles in the permutation $r_{n-2}r_n$ are coprime as they differ exactly by $1$, it follows that for any odd $a$ and even $b$ there is $s \in \mathbb{Z}^+$ such that $(r_{n}r_{n-2})^{s} (1\ 2)(r_{n-2}r_{n})^{s} = (a\ b)$. Thus, $H$ contains the Coxeter generators of $\Sym_n$, and therefore $H = \Sym_n$.

For the necessary statement, notice that if both $n,k$ have the same parity, then all generators of $H$ have indices that are all of the same parity. Therefore, by Observations $(3)$ and $(5)$ in the Introduction, $H<\Sym_n$. This completes the proof. 
\end{proof}

\end{document}
